\section{Exact proof of the certified input}\label{app:tail}

This appendix proves Lemma~\ref{lem:certified-input}. The positive-side
theta series has \(j\)-th derivative
\[
 2\pi n^2e^{5u/2-n^2x}R_j(n^2x),\qquad x=\pi e^{2u},
\]
where the polynomials \(R_j\) satisfy the recurrence in
\secref{sec:inputs}. Machin's identity and alternating series give
\(\pi>157/50\); hence \(R_0(x)=2x-3>0\).

\subsection{Cleared sign polynomials}

In addition to \(a_j\), define the normalized remaining modes
\[
 e_j(x)=\sum_{n\geq3}
 \frac{n^2e^{-(n^2-1)x}R_j(n^2x)}{R_0(x)},
 \qquad \widehat a_j=a_j+e_j.
\]
For an indeterminate raw jet \((\widehat a_0,\ldots,\widehat a_6)\), put
\[
 N_q=\widehat a_1^2-\widehat a_0\widehat a_2,
\]
\begin{align*}
N_{F_2}={}&-\widehat a_0^4\widehat a_2\widehat a_4
+\widehat a_0^4\widehat a_3^2
+\widehat a_0^3\widehat a_1^2\widehat a_4
-2\widehat a_0^3\widehat a_1\widehat a_2\widehat a_3
+2\widehat a_0^3\widehat a_2^3\\
&-3\widehat a_0^2\widehat a_1^2\widehat a_2^2
+3\widehat a_0\widehat a_1^4\widehat a_2-\widehat a_1^6,
\end{align*}
and
\[
 N_{\Cfour}=\det[\widehat a_{i+j}]_{i,j=0}^3.
\]
Direct logarithmic differentiation and the central-moment determinant give
\begin{equation}\label{eq:raw-signs}
 q=\frac{N_q}{\widehat a_0^2},\qquad
 F_2=\frac{N_{F_2}}{\widehat a_0^6},\qquad
 \Cfour=\frac{N_{\Cfour}}{\widehat a_0^4}.
\end{equation}
The thirteen-term cumulant expansion obtained from the last identity agrees
with \eqref{eq:C4-expanded}; this is an exact polynomial identity.

For orientation, retaining only the first mode gives
\[
q_1=\frac{4x(4x^2-12x+15)}{(2x-3)^2},
\]
\[
(F_2)_1=\frac{64x^3(64x^6-576x^5+2448x^4-6432x^3
+10044x^2-8100x+2295)}{(2x-3)^6},
\]
\[
(\Cfour)_1=\frac{49152x^6(16x^4-96x^3+360x^2-840x+945)}{(2x-3)^4}.
\]
After \(x=5+z\), every numerator has positive coefficients. The second
mode is retained exactly because it supplies the larger margin near the
origin.

\subsection{Exact two-mode margin}

Rational Taylor bounds for the exponential imply
\[
 0<\eta<\frac1{12000}\quad(x\geq\pi),
\]
and
\[
 \frac1{23000}<\eta<\frac1{12000}
 \quad\left(\frac{157}{50}\leq x\leq\frac{10}{3}\right).
\]
The close endpoints are entirely rational checks. Namely,
\[
 \sum_{k=0}^{16}\frac{(471/50)^k}{k!}>12000,\qquad
 \sum_{k=0}^{21}\frac{10^k}{k!}>22000,
\]
and the geometric bound for the tail after degree \(13\) gives
\[
 e^{10}<
 \sum_{k=0}^{13}\frac{10^k}{k!}
 \frac{10^{14}/14!}{1-10/15}
 =\frac{3480060751}{154791}<23000.
\]
These prove, respectively, the upper \(\eta\)-cap, the cap
\(e^{-10}<1/22000\), and the lower bound \(e^{-10}>1/23000\).
Likewise \(\sum_{k=0}^{21}15^k/k!>3{,}000{,}000\), which gives
\(\eta<1/3{,}000{,}000\) for \(x\geq5\).
We use the following elementary certificate lemma. If
\(p(s,t)=\sum p_{ij}s^it^j\) has bidegree \((d,e)\), its Bernstein
coefficients on the unit square are
\[
 \beta_{k\ell}=\sum_{i\leq k,\,j\leq\ell}
 p_{ij}\frac{\binom{k}{i}}{\binom{d}{i}}
       \frac{\binom{\ell}{j}}{\binom{e}{j}}.
\]
The power-to-Bernstein identity proves that \(\beta_{k\ell}>0\) for every
\(k,\ell\) implies \(p>0\) on the square.

Apply this lemma after affine rational changes of variables. For \(q\), one
rectangle ends at \(x=10/3\), followed by a positive base polynomial minus
exponentially decreasing corrections. For \(F_2\), dependent bounds on the
first rectangle are followed by a rectangle through \(x=5\) and a positive
half-strip expansion. For \(\Cfour\), the first rectangle is followed by two
negative odd-power corrections whose ratios to the positive base decrease.
The derivative of each ratio has a numerator with positive coefficients after
the indicated shift. The 361 resulting rational coefficient inequalities give
\begin{equation}\label{eq:two-mode-margin}
 q_{1,2}>10,\qquad (F_2)_{1,2}>1000,\qquad
 (\Cfour)_{1,2}>50{,}000{,}000.
\end{equation}
This is a finite algebraic proof on fixed domains.

Here is the precise connection between those coefficient checks and the
quantitative margins. Evaluate the raw polynomials at the two-mode jet and set
\[
 \mathcal Q=N_q-10a_0^2,
 \qquad \mathcal F=N_{F_2}-1000a_0^6,
 \qquad \mathcal C=N_{\Cfour}-50{,}000{,}000a_0^4.
\]
The rows labelled \(Q,F,C\) below certify the cleared numerators of
\(\mathcal Q,\mathcal F,\mathcal C\), respectively. Thus their positivity is
exactly, rather than merely qualitatively, the three inequalities in
\eqref{eq:two-mode-margin}. For a ratio row, write the shifted target as
\(B_0(x)+\sum_{j\geq1}\eta^jB_j(x)\), where \(B_0\) has positive shifted
coefficients. Whenever \(B_j\) is adverse, put \(D_j=-B_j\). Positivity of
\[
 3jD_jB_0-D_j'B_0+D_jB_0'
\]
after the same shift proves that
\(e^{-3jx}D_j(x)/B_0(x)\) decreases. The displayed endpoint ratio is the sum
of all such adverse ratios at the left endpoint; being less than one proves
the target positive on the entire half-line. Coefficientwise nonnegative
pieces only increase it.

For reference, the complete finite certificate is summarized below.  Write
\(\mathcal A=(2x-3)+4(8x-3)\eta\), which is positive on every listed domain.
Rows \(Q_0,F_0,F_1,C_0\) are Bernstein certificates after the indicated
affine change to the unit square; the remaining rows use coefficient
positivity after a half-line shift. In the two ``ratio'' rows, the listed
fraction is the sum of
the adverse coefficient ratios and is strictly less than one.
\begin{center}
\tiny
\begin{tabular}{@{}lllll@{}}
\toprule
ID & Variables and domain & Degree & Positive denominator & Exact lower datum \\
\midrule
\(Q_0\) & \(157/50\le x\le10/3,\ 0\le \eta\le1/12000\)
 & \((4,2)\) & \(\mathcal A^2\) & \(6613625187767/70312500000\) \\
\(Q_\infty\) & \(x=10/3+z,\ 0\le \eta\le1/12000\)
 & \((4,2)\) & \(\mathcal A^2\) & ratio \(47333/505500\) \\
\(F_0\) & \(157/50\le x\le10/3,\ 1/23000\le \eta\le1/12000\)
 & \((12,6)\) & \(\mathcal A^6\) & minimum \(>1610047\) \\
\(F_1\) & \(10/3\le x\le5,\ 0\le \eta\le1/22000\)
 & \((12,6)\) & \(\mathcal A^6\) & \(29124973000/19683\) \\
\(F_\infty\) & \(x=5+z,\ \eta=v/3000000,\ 0\leq v\leq1\)
 & \((12,6)\) & \(\mathcal A^6\) & \(432/19073486328125\) \\
\(C_0\) & \(157/50\le x\le10/3,\ 0\le \eta\le1/12000\)
 & \((14,4)\) & \(\mathcal A^4\) & minimum \(>9546672947\) \\
\(C_\infty\) & \(x=10/3+z,\ 0\le \eta\le1/12000\)
 & \((14,4)\) & \(\mathcal A^4\) & ratio \(65989337396066791/77165720550000000\) \\
\bottomrule
\end{tabular}
\end{center}
These seven checks comprise 361 exact rational inequalities.  The derivative
envelopes used below contribute 140 further Bernstein coefficients, of
bidegrees \((j+1,1)\), \(0\le j\le6\); their smallest coefficient is
\(44861/31250\).  Thus the compact domains and the differentiated tail
majorants belong to the same exact certificate.

\subsection{All later modes}

For direct reconstruction of the derivative bounds, put
\(S_j(t)=2^j(-1)^jR_j(t)\). The recurrence gives
\begin{align*}
S_0={}&2t-3,\\
S_1={}&8t^2-30t+15,\\
S_2={}&32t^3-224t^2+330t-75,\\
S_3={}&128t^4-1440t^3+4232t^2-3270t+375,\\
S_4={}&512t^5-8448t^4+41408t^3-68096t^2+30930t-1875,\\
S_5={}&2048t^6-46592t^5+343040t^4-976320t^3
       +1008968t^2-285870t+9375,\\
S_6={}&8192t^7-245760t^6+2536960t^5-11109120t^4
       +20633312t^3-14260064t^2+2610330t-46875.
\end{align*}

For \(t\geq27\), shifting \(t=27+z\) gives positive coefficients in
\((-1)^jR_j(t)\) and in
\(2^{j+1}t^{j+1}-(-1)^jR_j(t)\). Therefore
\begin{equation}\label{eq:R-envelope}
 0<(-1)^jR_j(t)\leq2^{j+1}t^{j+1},
 \qquad 0\leq j\leq6.
\end{equation}
The derivative numerator of the normalized \(n=3\) term is another positive
shifted polynomial, so its absolute value decreases for \(x\geq3\). For
\(n\geq4\), successive majorants have ratio at most
\[
 \left(\frac54\right)^{16}e^{-27}<10^{-9}.
\]
Rational Taylor lower bounds for \(e^{24},e^{27},e^{45}\) then yield, on
\(\pi\leq x\leq5\),
\begin{equation}\label{eq:core-tail}
(|e_0|,\ldots,|e_6|)<
(7\!\cdot\!10^{-9},4\!\cdot\!10^{-7},2\!\cdot\!10^{-4},
7\!\cdot\!10^{-4},3\!\cdot\!10^{-2},1.1,41).
\end{equation}
Indeed, the normalized \(n=3\) term decreases on \(x\geq3\), and
\(e^{24}>25{,}000{,}000{,}000\), so it is bounded by
\(3|R_j(27)|/25{,}000{,}000{,}000\). The sum of the \(n\geq4\) majorants is
at most
\[
 \frac{2^{j+2}3^j4^{2j+4}}{3\cdot10^{19}}.
\]
The data entering these two expressions are
\begin{center}
\small
\begin{tabular}{@{}ccl@{}}
\toprule
\(j\) & \( |R_j(27)| \) & chosen upper bound for \( |e_j| \)\\
\midrule
0 & \(51\) & \(7\cdot10^{-9}\)\\
1 & \(5037/2\) & \(4\cdot10^{-7}\)\\
2 & \(475395/4\) & \(2\cdot10^{-4}\)\\
3 & \(42678141/8\) & \(7\cdot10^{-4}\)\\
4 & \(3623251731/16\) & \(3\cdot10^{-2}\)\\
5 & \(288709329165/32\) & \(11/10\)\\
6 & \(21373315648611/64\) & \(41\)\\
\bottomrule
\end{tabular}
\end{center}
For each row the later-mode expression is less than one thousandth of the
\(n=3\) expression, and their sum is strictly below the chosen bound. This
derives every constant in \eqref{eq:core-tail} from the printed recurrence and
three rational exponential inequalities.
The same Bernstein lemma gives
\[
(|a_0|,\ldots,|a_6|)<(2,5,20,200,2000,60000,1000000).
\]
For any one of the printed raw polynomials
\(N=\sum_\alpha c_\alpha X^\alpha\), the perturbation estimate used here is
the explicit finite sum
\begin{equation}\label{eq:perturbation-formula}
 |N(a+e)-N(a)|
 \leq\sum_\alpha |c_\alpha|
 \left\{\prod_{j=0}^6(A_j+E_j)^{\alpha_j}
              -\prod_{j=0}^6A_j^{\alpha_j}\right\},
\end{equation}
whenever \(|a_j|\leq A_j\) and \(|e_j|\leq E_j\). Substitution of the two
displayed seven-tuples into \eqref{eq:perturbation-formula} gives
\begin{equation}\label{eq:core-perturbation}
 |\Delta N_q|<0.000405,\quad
 |\Delta N_{F_2}|<37.959,\quad
 |\Delta N_{\Cfour}|<31{,}549{,}532.
\end{equation}

For \(x\geq5\), \eqref{eq:R-envelope} gives
\[
 |a_j|\leq2^{j+2}x^j,\qquad
 |e_j|\leq2^{j+2}3^{2j+4}x^je^{-8x}.
\]
The raw weights of \(N_q,N_{F_2},N_{\Cfour}\) are respectively \(2,6,12\).
Thus every perturbation majorant is a positive multiple of
\(x^we^{-8x}\), or a faster-decaying term, with \(w\leq12\); it decreases
for \(x\geq5\). In \eqref{eq:perturbation-formula} take
\[
 A_j=2^{j+2}5^j,\qquad
 E_j=\frac{2^{j+1}3^{2j+4}5^j}{10^{17}}.
\]
The latter follows from the displayed tail bound and
\(\sum_{k=0}^{47}40^k/k!>2\cdot10^{17}\). Direct substitution gives
\begin{equation}\label{eq:outer-perturbation}
 |\Delta N_q|<6.49\cdot10^{-11},\quad
 |\Delta N_{F_2}|<0.03,\quad
 |\Delta N_{\Cfour}|<418{,}287.
\end{equation}
To compare like quantities, note that
\[
 a_0=1+4\eta\frac{R_0(4x)}{R_0(x)}>1,
\]
and, by homogeneity,
\[
 N_{q,1,2}=q_{1,2}a_0^2,\qquad
 N_{F_2,1,2}=(F_2)_{1,2}a_0^6,\qquad
 N_{\Cfour,1,2}=(\Cfour)_{1,2}a_0^4.
\]
Consequently \eqref{eq:two-mode-margin} supplies the same respective lower
margins \(10\), \(1000\), and \(50{,}000{,}000\) for the cleared
numerators.  Both perturbation bounds
\eqref{eq:core-perturbation} and \eqref{eq:outer-perturbation} are strictly
smaller. Their scale relative to the corresponding cleared-numerator margin
is recorded here:
\begin{center}
\small
\begin{tabular}{@{}lcc@{}}
\toprule
target & surviving fraction on \(\pi\leq x\leq5\)
       & surviving fraction on \(x\geq5\)\\
\midrule
\(q\) & \(>0.9999595\) & \(>0.99999999999351\)\\
\(F_2\) & \(>0.962041\) & \(>0.99997\)\\
\(\Cfour\) & \(>0.36900936\) & \(>0.99163426\)\\
\bottomrule
\end{tabular}
\end{center}
Thus even the least relative margin, the compact \(\Cfour\) bound, retains
more than \(36.9\) percent of its certified two-mode numerator margin.
Finally \(\widehat a_0=a_0+e_0>0\), so \eqref{eq:raw-signs} proves all three
signs for \(u\geq0\). Evenness proves Lemma~\ref{lem:certified-input}.
